% from ega3-1-fr.tex
\oldpage[0\textsubscript{III}]{6}
$g(Y)$ est donc une application $h_X(Y)\to F(Y)$ telle que pour tout morphisme
$u:Y\to Y'$ dans $C$, le diagramme
\[
  \xymatrix{
    h_X(Y')\ar[r]^{h_X(u)}\ar[d]_{g(Y')} &
    h_X(Y)\ar[d]^{g(Y)}\\
    F(Y')\ar[r]_{F(u)} & F(Y)
  }
  \tag{8.1.3.1}\label{0.8.1.3.1-fr}
\]
soit commutatif. En particulier, on a une application
$g(X):h_X(X)=\operatorname{Hom}(X,X)\to F(X)$, d'où un élément
\[
  \alpha(g)=(g(X))(1_X)\in F(X)
  \tag{8.1.3.2}
\]
et par suite une application canonique
\[
  \alpha:\operatorname{Hom}(h_X,F)\to F(X).
  \tag{8.1.3.3}
\]

Inversement, considérons un élément $\xi\in F(X)$~; pour tout morphisme
$v:Y\to X$ dans $C$, $F(v)$ est une application $F(X)\to F(Y)$~;
considérons l'application
\[
  v\mapsto(F(v))(\xi)
  \tag{8.1.3.4}
\]
de $h_X(Y)$ dans $F(Y)$~; si on désigne par $(\beta(\xi))(Y)$ cette
application,
\[
  \beta(\xi):h_X\to F
  \tag{8.1.3.5}
\]
est un \emph{morphisme fonctoriel}, car on a pour tout morphisme
$u:Y\to Y'$ dans $C$,
$(F(vu))(\xi)=(F(v)\circ F(u))(\xi)$, ce qui vérifie la commutativité de
(\hyperref[0.8.1.3.1-fr]{8.1.3.1}) pour $g=\beta(\xi)$. On a ainsi défini
une application canonique
\[
  \beta:F(X)\to\operatorname{Hom}(h_X,F).
  \tag{8.1.3.6}
\]
\end{env}

\begin{proposition}[8.1.4]
\label{0.8.1.4-fr}
Les applications $\alpha$ et $\beta$ sont des bijections réciproques l'une de
l'autre.
\end{proposition}

Calculons $\alpha(\beta(\xi))$ pour $\xi\in F(X)$~; pour tout $Y\in C$,
$(\beta(\xi))(Y)$ est l'application
$g_1(Y):v\mapsto(F(v))(\xi)$ de $h_X(Y)$ dans $F(Y)$. On a donc
\[
  \alpha(\beta(\xi))=(g_1(X))(1_X)=(F(1_X))(\xi)
  =1_{F(X)}(\xi)=\xi.
\]
Calculons maintenant $\beta(\alpha(g))$ pour
$g\in\operatorname{Hom}(h_X,F)$~; pour tout $Y\in C$,
$(\beta(\alpha(g)))(Y)$ est l'application
$v\mapsto(F(v))((g(X))(1_X))$~; en vertu de la commutativité de
(\hyperref[0.8.1.3.1-fr]{8.1.3.1}), cette application n'est autre que
$v\mapsto(g(Y))((h_X(v))(1_X))=(g(Y))(v)$ par définition de $h_X(v)$,
autrement dit, elle est égale à $g(Y)$, ce qui démontre la proposition.

\begin{env}[8.1.5]
\label{0.8.1.5-fr}
Rappelons qu'une \emph{sous-catégorie} $C'$ d'une catégorie $C$ est définie
par la condition que ses objets soient des objets de $C$, et que si $X'$,
$Y'$ sont deux objets de $C'$, l'ensemble
$\operatorname{Hom}_{C'}(X',Y')$ des morphismes $X'\to Y'$ \emph{dans
$C'$} est une partie de l'ensemble $\operatorname{Hom}_C(X',Y')$ des
morphismes $X'\to Y'$ \emph{dans $C$}, l'application canonique de
<<~composition des morphismes~>>
\[
  \operatorname{Hom}_{C'}(X',Y')\times\operatorname{Hom}_{C'}(Y',Z')
  \to\operatorname{Hom}_{C'}(X',Z')
\]
